\documentclass[10pt,a4paper]{amsart}
\usepackage[margin=19mm]{geometry}
\usepackage{amsmath,amssymb,mathtools}
\usepackage[scr=rsfs,cal=euler]{mathalfa}
\usepackage{xfrac}
\usepackage[unicode,hidelinks,bookmarks=false]{hyperref}
\newcommand{\ie}{i.e.\ }
\newcommand{\cf}{cf.\ }
\newcommand{\resp}{respectively\ }
\newcommand{\Subsection}{\subsection}
\providecommand{\nobreakdash}{\nobreak\hbox{-}}
\setlength{\emergencystretch}{2em}
\multlinegap0pt
\usepackage{bm}
\usepackage{bbm}
\usepackage{dsfont}
\usepackage{xspace}
\usepackage{cancel}
\usepackage{mathtools}
\usepackage{amsthm}
\makeatletter
\@ifundefined{lemma}{\newtheorem{lemma}{Lemma}}{}
\@ifundefined{proposition}{\newtheorem{proposition}{Proposition}}{}
\makeatother
\title[Strategic Risk Reduction: Self-Protection and Self-Insurance]{Strategic Risk Reduction: Self-Protection and Self-Insurance}
\author{Wing Fung Chong}
\date{Source: arXiv:2606.30363v2 (CC BY 4.0)}
\begin{document}
\maketitle

\noindent\emph{Source and licence.} Strategic Risk Reduction: Self-Protection and Self-Insurance. Version arXiv:2606.30363v2 (CC BY 4.0), distributed under the Creative Commons Attribution 4.0 International licence, as recorded in the arXiv licence metadata for this exact version and shown on the abstract page https://arxiv.org/abs/2606.30363v2. Full rights notice: licenses/arxiv-2606.30363-CC-BY-4.0.txt.\par\smallskip

\noindent\textit{Editorial scope.} Counted unit: the Proposition whose source label is \texttt{prop:tvar\_second\_result} in the article (source0.tex lines 622--629, source label \texttt{prop:tvar\_second\_result}), delivered together with the article's own complete original proof of it (source0.tex lines 631--650, one \texttt{proof} environment of 20 lines). No mathematics is written, repaired or re-derived: statement and proof are the article's own text line for line. Required context retained uncounted: eq:TVaR\_problem<= (lines 417--420); eq:gradient\_property (lines 534--543). Every definition, display and label of those lines is retained unchanged. Omitted from this package and not redistributed: 1--416, 421--533, 544--621, 630--630, 651--1475. Nothing inside a retained excerpt is omitted or altered: every retained body is delivered exactly as the article's lines are written, so no omission receipt of any kind accompanies this package. Citations: the retained excerpts carry no citation command at all, so no bibliography entry of the article is transcribed and no reference is redistributed. Reference adaptation: none was needed and none was made; every reference inside the delivered unit resolves inside it, so no unresolved reference or citation is produced. Printed slips kept literal: no printed word, symbol, hypothesis or number of the counted statement or of its proof was altered. Layout and class: the article's own class is \texttt{article} with packages including abstract, amsfonts, amsmath, bm, amssymb, amsthm, appendix, array, authblk, bbm, caption, dsfont, setspace, enumitem; the wrapper is the lane's pinned \texttt{amsart} editorial wrapper at 10pt on A4, which restates the theorem-like environments the retained text needs and adds the article's own macro definitions that the retained text needs (none), with extra wrapper packages (bm, bbm, dsfont, xspace, cancel, mathtools). The delivered numbering is the unit's own. Assets: no retained span contains a figure environment, a graphics-inclusion command, a PDF-page inclusion, a table or a TikZ picture; no image file is copied into this package, the package declares no asset dependency and no image file is redistributed with it. Licence: version arXiv:2606.30363v2 of this article is distributed under the Creative Commons Attribution 4.0 International licence, as recorded in the arXiv licence metadata for this exact version and shown on the abstract page https://arxiv.org/abs/2606.30363v2. A full rights notice is delivered at licenses/arxiv-2606.30363-CC-BY-4.0.txt. Characters: every retained excerpt is pure ASCII, so the delivered engine, XeLaTeX with its default Latin Modern text font, reports no missing character.
\begin{equation}
\inf_{\left(p,l\right)\in\Theta^{\left(\text{L}\right)}}G\left(p,l\right)=\inf_{\left(p,l\right)\in\Theta^{\left(\text{L}\right)}}\pi\left(p,l\right)+\frac{pl}{1-\alpha},
\label{eq:TVaR_problem<=}
\end{equation}

\begin{lemma}\label{eq:gradient_property}
Assume that $\Theta^{\left(\text{L}\right)}\neq\emptyset$ and $G_{pl}>0$ on $\Theta^{\left(\text{L}\right)}$. Then the following statements hold.
\begin{enumerate}
\item[(i)] For any $\left(p,l\right)\in\Theta_{\text{F}}^{\left(\leq\right)}$, and for any $\left(\tilde{p},\tilde{l}\right)\in\Theta^{\left(\text{L}\right)}\setminus\left\{\left(p,l\right)\right\}$ with $\tilde{p}\leq p$ and $\tilde{l}\leq l$, we have $\left(\tilde{p},\tilde{l}\right)\in\Theta_{\text{F}}^{\left(<\right)}$.
\item[(ii)] For any $\left(p,l\right)\in\Theta_{\text{F}}^{\left(\geq\right)}$, and for any $\left(\tilde{p},\tilde{l}\right)\in\Theta^{\left(\text{L}\right)}\setminus\left\{\left(p,l\right)\right\}$ with $\tilde{p}\geq p$ and $\tilde{l}\geq l$, we have $\left(\tilde{p},\tilde{l}\right)\in\Theta_{\text{F}}^{\left(>\right)}$.
\item[(iii)] For any $\left(p,l\right)\in\Theta_{\text{S}}^{\left(\leq\right)}$, and for any $\left(\tilde{p},\tilde{l}\right)\in\Theta^{\left(\text{L}\right)}\setminus\left\{\left(p,l\right)\right\}$ with $\tilde{p}\leq p$ and $\tilde{l}\leq l$, we have $\left(\tilde{p},\tilde{l}\right)\in\Theta_{\text{S}}^{\left(<\right)}$.
\item[(iv)] For any $\left(p,l\right)\in\Theta_{\text{S}}^{\left(\geq\right)}$, and for any $\left(\tilde{p},\tilde{l}\right)\in\Theta^{\left(\text{L}\right)}\setminus\left\{\left(p,l\right)\right\}$ with $\tilde{p}\geq p$ and $\tilde{l}\geq l$, we have $\left(\tilde{p},\tilde{l}\right)\in\Theta_{\text{S}}^{\left(>\right)}$.
\end{enumerate}
Consequently, if $\mathcal{I}_{\text{F}}=\Theta_{\text{F}}^{\left(=\right)}$ or $\mathcal{I}_{\text{S}}=\Theta_{\text{S}}^{\left(=\right)}$ is non-empty and can be represented as a graph, then it is downward-sloping. When such an isoquant cuts through the rectangle $\Theta^{\left(\text{L}\right)}$, it separates the rectangle into lower-left and upper-right sign regions.
\end{lemma}

\begin{proposition}\label{prop:tvar_second_result}
Assume that $\Theta^{\left(\text{L}\right)}\neq\emptyset$ and $G_{pl}>0$ on $\Theta^{\left(\text{L}\right)}$. Suppose that $G_p\left(\underline{p},\underline{l}\right)<0<G_p\left(\overline{p}_{\alpha},\overline{l}\right)$ and $G_l\left(\underline{p},\underline{l}\right)<0<G_l\left(\overline{p}_{\alpha},\overline{l}\right)$, but $\Theta_{\text{F,S}}^{\left(=,=\right)}=\emptyset$. Then the following statements hold.

\begin{enumerate}
\item[(i)] If $\Theta_{\text{F}}^{\left(=\right)}\subseteq\Theta_{\text{S}}^{\left(<\right)}$ or, equivalently, $\Theta_{\text{S}}^{\left(=\right)}\subseteq\Theta_{\text{F}}^{\left(>\right)}$, then the optimal solution of the left-region problem \eqref{eq:TVaR_problem<=} is ${\bm\theta}^{*\left(\text{L}\right)}={\bm\theta}^{\text{UL}}=\left(p^*\left(\overline{l}\right),l^*\left(\underline{p}\right)\right)$.
\item[(ii)] If $\Theta_{\text{F}}^{\left(=\right)}\subseteq\Theta_{\text{S}}^{\left(>\right)}$ or, equivalently, $\Theta_{\text{S}}^{\left(=\right)}\subseteq\Theta_{\text{F}}^{\left(<\right)}$, then the optimal solution of the left-region problem \eqref{eq:TVaR_problem<=} is ${\bm\theta}^{*\left(\text{L}\right)}={\bm\theta}^{\text{LR}}=\left(p^*\left(\underline{l}\right),l^*\left(\overline{p}_{\alpha}\right)\right)$.
\end{enumerate}
\end{proposition}

\begin{proof}
We prove (i). Suppose that $\Theta_{\text{F}}^{\left(=\right)}\subseteq\Theta_{\text{S}}^{\left(<\right)}$. Equivalently, $\Theta_{\text{S}}^{\left(=\right)}\subseteq\Theta_{\text{F}}^{\left(>\right)}$. Under this configuration, the frequency marginal-balance isoquant lies entirely in the region where the severity derivative is negative, while the severity marginal-balance isoquant lies entirely in the region where the frequency derivative is positive. Therefore, the feasible rectangle is decomposed as
\begin{equation*}
\Theta^{\left(\text{L}\right)}=\Theta_{\text{F,S}}^{\left(<,<\right)}\cup\Theta_{\text{F,S}}^{\left(=,<\right)}\cup\Theta_{\text{F,S}}^{\left(>,<\right)}\cup\Theta_{\text{F,S}}^{\left(>,=\right)}\cup\Theta_{\text{F,S}}^{\left(>,>\right)}.
\end{equation*}
Indeed, the joint sign regions $\Theta_{\text{F,S}}^{\left(<,=\right)}$, $\Theta_{\text{F,S}}^{\left(<,>\right)}$, $\Theta_{\text{F,S}}^{\left(=,=\right)}$, $\Theta_{\text{F,S}}^{\left(=,>\right)}$ are empty under the assumed configuration.

We first show that no optimal solution can lie in the two regions $\Theta_{\text{F,S}}^{\left(<,<\right)}$ and $\Theta_{\text{F,S}}^{\left(>,>\right)}$. Let $\left(p,l\right)\in\Theta_{\text{F,S}}^{\left(<,<\right)}$. Then $G_p\left(p,l\right)<0$ and $G_l\left(p,l\right)<0$. By Lemma \ref{eq:gradient_property}, there exists a feasible northeast movement from $\left(p,l\right)$ along which the objective strictly decreases until the path reaches the boundary of $\Theta_{\text{F,S}}^{\left(=,<\right)}\cup\Theta_{\text{F,S}}^{\left(>,<\right)}\cup\Theta_{\text{F,S}}^{\left(>,=\right)}$. Along such a movement, the directional derivative of $G$ is negative. Hence every strategy in $\Theta_{\text{F,S}}^{\left(<,<\right)}$ is strictly dominated by some strategy in $\Theta_{\text{F,S}}^{\left(=,<\right)}\cup\Theta_{\text{F,S}}^{\left(>,<\right)}\cup\Theta_{\text{F,S}}^{\left(>,=\right)}$. Thus no strategy in $\Theta_{\text{F,S}}^{\left(<,<\right)}$ can be optimal. Similar arguments show that every strategy in $\Theta_{\text{F,S}}^{\left(>,>\right)}$ is strictly dominated by some strategy in $\Theta_{\text{F,S}}^{\left(=,<\right)}\cup\Theta_{\text{F,S}}^{\left(>,<\right)}\cup\Theta_{\text{F,S}}^{\left(>,=\right)}$, and thus no strategy in $\Theta_{\text{F,S}}^{\left(>,>\right)}$ can be optimal.

Next, we show that ${\bm\theta}^{\text{UL}}=\left(p^{\text{UL}},l^{\text{UL}}\right)\in\Theta_{\text{F,S}}^{\left(=,<\right)}\cup\Theta_{\text{F,S}}^{\left(>,<\right)}\cup\Theta_{\text{F,S}}^{\left(>,=\right)}$, and that every strategy in this set lies southeast of ${\bm\theta}^{\text{UL}}$. Under the assumed configuration, there are the following three cases.
\begin{enumerate}
\item[(I)] Suppose that $\left(\underline{p},\overline{l}\right)\in\Theta_{\text{F}}^{\left(\leq\right)}$. Then, under the present ordering, $\left(\underline{p},\overline{l}\right)\in\Theta_{\text{S}}^{\left(<\right)}$. In this case, ${\bm\theta}^{\text{UL}}=\left(p^{\text{UL}},l^{\text{UL}}\right)=\left(p^*\left(\overline{l}\right),\overline{l}\right)$. Moreover, $p^*\left(\overline{l}\right)\in\left[\underline{p},\overline{p}_{\alpha}\right)$ is the constrained marginal-balance strategy on the upper boundary; if it is interior, then it satisfies $G_p\left(p^*\left(\overline{l}\right),\overline{l}\right)=0$. Thus ${\bm\theta}^{\text{UL}}\in\Theta_{\text{F,S}}^{\left(=,<\right)}\subseteq\Theta_{\text{F,S}}^{\left(=,<\right)}\cup\Theta_{\text{F,S}}^{\left(>,<\right)}\cup\Theta_{\text{F,S}}^{\left(>,=\right)}$. Let $\left(p,l\right)\in\Theta_{\text{F,S}}^{\left(=,<\right)}\cup\Theta_{\text{F,S}}^{\left(>,<\right)}\cup\Theta_{\text{F,S}}^{\left(>,=\right)}$. Clearly, $l\leq\overline{l}=l^{\text{UL}}$. If $p^{\text{UL}}=\underline{p}$, then $p^{\text{UL}}\leq p$. If $p^{\text{UL}}>\underline{p}$, suppose, to the contrary, that $p<p^{\text{UL}}$. Since ${\bm\theta}^{\text{UL}}\in\Theta_{\text{F,S}}^{\left(=,<\right)}\subseteq\Theta_{\text{F}}^{\left(=\right)}$, as well as $p<p^{\text{UL}}$ and $l\leq l^{\text{UL}}$, Lemma \ref{eq:gradient_property} implies that $\left(p,l\right)\in\Theta_{\text{F}}^{\left(<\right)}$, which contradicts $\left(p,l\right)\in\Theta_{\text{F,S}}^{\left(=,<\right)}\cup\Theta_{\text{F,S}}^{\left(>,<\right)}\cup\Theta_{\text{F,S}}^{\left(>,=\right)}$. Hence $p^{\text{UL}}\leq p$.
\item[(II)] When $\left(\underline{p},\overline{l}\right)\in\Theta_{\text{F,S}}^{\left(>,<\right)}$, ${\bm\theta}^{\text{UL}}=\left(p^{\text{UL}},l^{\text{UL}}\right)=\left(\underline{p},\overline{l}\right)$, and hence ${\bm\theta}^{\text{UL}}\in\Theta_{\text{F,S}}^{\left(>,<\right)}$. For any $\left(p,l\right)\in\Theta_{\text{F,S}}^{\left(=,<\right)}\cup\Theta_{\text{F,S}}^{\left(>,<\right)}\cup\Theta_{\text{F,S}}^{\left(>,=\right)}$, we have immediately $p^{\text{UL}}=\underline{p}\leq p$ and $l\leq\overline{l}=l^{\text{UL}}$.
\item[(III)] Suppose that $\left(\underline{p},\overline{l}\right)\in\Theta_{\text{S}}^{\left(\geq\right)}$. Then, under the present ordering, $\left(\underline{p},\overline{l}\right)\in\Theta_{\text{F}}^{\left(>\right)}$. In this case, ${\bm\theta}^{\text{UL}}=\left(p^{\text{UL}},l^{\text{UL}}\right)=\left(\underline{p},l^*\left(\underline{p}\right)\right)$. Moreover, $l^*\left(\underline{p}\right)\in\left(\underline{l},\overline{l}\right]$ is the constrained marginal-balance strategy on the left boundary; if it is interior, then it satisfies $G_l\left(\underline{p},l^*\left(\underline{p}\right)\right)=0$. Thus ${\bm\theta}^{\text{UL}}\in\Theta_{\text{F,S}}^{\left(>,=\right)}\subseteq\Theta_{\text{F,S}}^{\left(=,<\right)}\cup\Theta_{\text{F,S}}^{\left(>,<\right)}\cup\Theta_{\text{F,S}}^{\left(>,=\right)}$. Let $\left(p,l\right)\in\Theta_{\text{F,S}}^{\left(=,<\right)}\cup\Theta_{\text{F,S}}^{\left(>,<\right)}\cup\Theta_{\text{F,S}}^{\left(>,=\right)}$. Clearly, $p^{\text{UL}}=\underline{p}\leq p$. If $l^{\text{UL}}=\overline{l}$, then $l\leq l^{\text{UL}}$. If $l^{\text{UL}}<\overline{l}$, suppose, to the contrary, that $l>l^{\text{UL}}$. Since ${\bm\theta}^{\text{UL}}\in\Theta_{\text{F,S}}^{\left(>,=\right)}\subseteq\Theta_{\text{S}}^{\left(=\right)}$, as well as $p^{\text{UL}}\leq p$ and $l^{\text{UL}}<l$, Lemma \ref{eq:gradient_property} implies that $\left(p,l\right)\in\Theta_{\text{S}}^{\left(>\right)}$, which contradicts $\left(p,l\right)\in\Theta_{\text{F,S}}^{\left(=,<\right)}\cup\Theta_{\text{F,S}}^{\left(>,<\right)}\cup\Theta_{\text{F,S}}^{\left(>,=\right)}$. Hence $l\leq l^{\text{UL}}$.
\end{enumerate}

Combining the three cases, for every $\left(\tilde{p},\tilde{l}\right)\in\left(\Theta_{\text{F,S}}^{\left(=,<\right)}\cup\Theta_{\text{F,S}}^{\left(>,<\right)}\cup\Theta_{\text{F,S}}^{\left(>,=\right)}\right)\setminus\left\{{\bm\theta}^{\text{UL}}\right\}$, we have $p^{\text{UL}}\leq\tilde{p}$ and $\tilde{l}\leq l^{\text{UL}}$, with at least one inequality strict. Moreover, throughout $\Theta_{\text{F,S}}^{\left(=,<\right)}\cup\Theta_{\text{F,S}}^{\left(>,<\right)}\cup\Theta_{\text{F,S}}^{\left(>,=\right)}$, we have $G_p\geq 0$ and $G_l\leq 0$. By the separation geometry in Lemma \ref{eq:gradient_property}, there exists a monotone path from $\left(\tilde{p},\tilde{l}\right)$ to ${\bm\theta}^{\text{UL}}$, contained in $\Theta_{\text{F,S}}^{\left(=,<\right)}\cup\Theta_{\text{F,S}}^{\left(>,<\right)}\cup\Theta_{\text{F,S}}^{\left(>,=\right)}$, along which $p$ decreases and $l$ increases, with at least one coordinate changing strictly. Along this path, $\text{d}G=G_p\text{d}p+G_l\text{d}l\leq 0$, with strict inequality on a set of positive path length. Hence, by the gradient theorem, $G\left({\bm\theta}^{\text{UL}}\right)<G\left(\tilde{p},\tilde{l}\right)$. Consequently, ${\bm\theta}^{*\left(\text{L}\right)}={\bm\theta}^{\text{UL}}=\left(p^*\left(\overline{l}\right),l^*\left(\underline{p}\right)\right)$. This proves (i).

The proof of (ii) is analogous, with the roles of the upper-left and lower-right candidates reversed. Under the ordering $\Theta_{\text{F}}^{\left(=\right)}\subseteq\Theta_{\text{S}}^{\left(>\right)}$ or, equivalently, $\Theta_{\text{S}}^{\left(=\right)}\subseteq\Theta_{\text{F}}^{\left(<\right)}$, the joint sign regions that an optimal solution can lie in, after eliminating $\Theta_{\text{F,S}}^{\left(<,<\right)}$ and $\Theta_{\text{F,S}}^{\left(>,>\right)}$, are $\Theta_{\text{F,S}}^{\left(<,=\right)}\cup\Theta_{\text{F,S}}^{\left(<,>\right)}\cup\Theta_{\text{F,S}}^{\left(=,>\right)}$. The same argument shows that every strategy in this set lies northwest of ${\bm\theta}^{\text{LR}}$, and that moving southeast toward ${\bm\theta}^{\text{LR}}$ decreases $G$, with strict decrease unless the starting point is ${\bm\theta}^{\text{LR}}$ itself in the set. Hence ${\bm\theta}^{*\left(\text{L}\right)}={\bm\theta}^{\text{LR}}=\left(p^*\left(\underline{l}\right),l^*\left(\overline{p}_{\alpha}\right)\right)$. This proves (ii).
\end{proof}

\end{document}
